\chapter{A CASA PRIMA CHE FACCIA BUIO}\label{a-casa-prima-che-faccia-buio}

\hfill\break

E così trascorsi qualche giorno in un ospedale nell'area di Miami, a
ristabilirmi dalle ferite lievi, a descrivere gli avvenimenti agli
investigatori federali e ad affrontare la morte di Wun. Fu durante
questo periodo che nacque la mia determinazione di lasciare la Perielio
e aprire un ambulatorio privato.

Ma decisi di non comunicare la mia intenzione fin dopo il lancio dei
replicatori. Non volevo turbare Jason in un momento critico.

\separatorefiori

In confronto allo sforzo di terraformazione degli anni precedenti, il
lancio dei replicatori fu deludente. I suoi risultati, se ce ne fossero
stati, sarebbero stati più grandi e ingegnosi; ma la sua efficienza -
nient'altro che una manciata di missili senza alcuna sincronizzazione
intelligente - non fu niente di spettacolare.

Il Presidente Lomax aveva mantenuto la faccenda entro i confini
nazionali. Con un'iniziativa che aveva fatto infuriare l'U.E., i cinesi,
i russi e gli indiani, Lomax aveva rifiutato di condividere la
tecnologia dei replicatori al di fuori degli ambienti di competenza
della NASA e della Perielio, e aveva fatto eliminare dagli archivi
marziani desegretati tutti i passaggi rilevanti. I ``microbi
artificiali'' (nel gergo di Lomax) erano una tecnologia ``ad alto
rischio''. Potevano essere ``usati come arma''. (Questo era vero, come
anche Wun aveva riconosciuto.) Gli Stati Uniti erano pertanto obbligati
ad assumere il ``controllo detentivo'' delle informazioni al fine di
prevenire ``la proliferazione delle nanotecnologie e una nuova e mortale
corsa alle armi''.

L'Unione Europea aveva espresso il suo disappunto e le Nazioni Unite
stavano convocando una commissione d'indagine, ma in un mondo in cui i
conflitti divampavano in quattro continenti, le argomentazioni di Lomax
avevano un peso considerevole. (Anche se, come Wun avrebbe replicato, i
marziani erano riusciti a vivere con la stessa tecnologia per centinaia
di anni - e i marziani non erano né più né meno umani dei loro antenati
terrestri.)

Per tutte queste ragioni, la data del lancio di fine estate a Canaveral
attirò un numero minimo di spettatori e l'attenzione pressoché sporadica
dei media. Wun Ngo Wen era morto, dopotutto, e le agenzie di stampa
avevano dato fondo alle loro risorse per coprire il suo assassinio.
Ormai i quattro grossi razzi Delta collocati nelle loro rampe di lancio
in mare aperto sembravano poco più di una nota a piè di pagina della sua
funzione commemorativa o, peggio, una replica: lanci di semi
riconfigurati per un'epoca di aspettative ridotte.

Ma anche se si trattava di uno spettacolo secondario, era tuttavia uno
spettacolo. Lomax per l'occasione giunse in aereo. E.D. Lawton aveva
accettato l'invito di cortesia e durante tutto il tempo era intenzionato
a comportarsi bene. E così, il mattino del giorno stabilito, andai con
Jason alle tribune dei VIP sulla costa orientale di Cape Canaveral.

Guardammo verso il mare. Le vecchie piattaforme al largo, ancora
funzionanti ma un po' ossidate per la ruggine provocata dall'acqua
salata, erano state costruite per sostenere i trasporti più pesanti
dell'era dei lanci di semi. I nuovissimi Delta, al confronto, sembravano
rimpiccioliti. Non che potessimo vedere bene i dettagli da quella
distanza, solo quattro colonne bianche all'orizzonte nebbioso
dell'oceano d'estate, oltre alla struttura di altre piattaforme di
lancio inutilizzate, connettori, mezzi ausiliari e navi di appoggio
ancorate a un perimetro di sicurezza. Era una mattina d'estate limpida e
calda. Il vento soffiava a raffiche - non tanto forte da annullare il
lancio ma abbastanza per sferzare la bandiera e scompigliare
l'acconciatura del Presidente Lomax mentre saliva sul podio per
rivolgersi ai dignitari e ai membri della stampa riuniti.

Pronunciò un discorso, breve per fortuna. Menzionò il lascito di Wun Ngo
Wen e la sua fiducia nel fatto che la rete di replicatori che era in
procinto di essere piantata ai confini gelidi del sistema solare avrebbe
presto fatto luce sulla natura e lo scopo dello Spin. Disse cose
magnifiche a proposito dell'umanità che lascia il suo segno nel cosmo.

«Intende dire la galassia,» bisbigliò Jason, «non il cosmo. E\ldots{}
lasciare il nostro \emph{segno}? Come un cane che piscia su un idrante?
Qualcuno dovrebbe davvero rivedere questi discorsi.»

Poi Lomax citò una poesia di un poeta russo del diciannovesimo secolo di
nome F.I. Tjutcev, il quale non avrebbe potuto immaginare lo Spin ma
scrisse come se l'avesse fatto:

``\emph{Svanito come una visione è il mondo esterno}

\emph{e l'uomo, orfano senza casa, deve affrontare}

\emph{indifeso, nudo e solo},

\emph{l'oscurità dell'incommensurabile spazio}.

\emph{Tutta la vita e lo splendore sembrano un antico sogno},

\emph{mentre nell'essenza della notte},

\emph{rivelata, estranea, egli ora avverte}

\emph{qualcosa di fatidico che è suo di diritto}.''

A quel punto Lomax abbandonò il palco e, dopo il compito prosaico del
conteggio alla rovescia, il primo razzo si levò in una colonna di fuoco
nel cosmo che si andò rivelando al di là del cielo. Qualcosa di
fatidico. Nostro di diritto.

Mentre tutti gli altri guardavano in alto, Jason chiuse gli occhi e si
strinse le mani sul grembo.

\separatorefiori

Ci spostammo nella sala di ricevimento insieme al resto degli invitati,
in attesa di un giro di interviste da parte della stampa. (Jason era in
programma per un colloquio di venti minuti sul canale dei notiziari, io
per dieci minuti. Io ero ``il dottore che tentò di salvare la vita di
Wun Ngo Wen'' anche se non avevo fatto altro che spegnergli la scarpa in
fiamme e tirarlo via dalla linea di fuoco dopo che era caduto. Un rapido
controllo di base - vie aeree, respirazione, circolazione - mi aveva
abbondantemente fatto capire che non potevo aiutarlo e che sarebbe stato
più saggio tenere semplicemente giù la testa fino all'arrivo dei
soccorsi. Era ciò che continuavo a raccontare ai cronisti, finché non
smisero di chiedermelo.)

Il Presidente Lomax attraversò la sala stringendo mani prima di essere
spinto via ancora una volta dalle sue guardie del corpo. Poi E.D. bloccò
Jason e me al tavolo del buffet.

«Credo che tu abbia ottenuto ciò che volevi» disse, indirizzando il
commento a Jason ma guardando me. «Adesso non può essere annullato.»

«In questo caso,» lo provocò Jason, «forse non vale la pena discuterne.»

Io e Wun avevamo ritenuto importante tenere Jase sotto osservazione nei
mesi successivi al suo trattamento. Si era sottoposto a una batteria di
test neurologici tra cui un'altra serie di risonanze magnetiche segrete.
Nessuno dei test aveva rivelato deficienze, e gli unici cambiamenti
fisiologici evidenti erano legati al recupero dalla SMA. In altre
parole, un certificato di buona salute. Migliore di quello che un tempo
avrei potuto immaginare possibile.

Però sembrava leggermente diverso. Avevo chiesto a Wun se tutti i Quarti
subissero cambiamenti psicologici.

«In un certo senso,» aveva risposto, «sì.» Ci si attendeva che i Quarti
si comportassero in modo differente dopo il loro trattamento, ma c'era
una certa sottigliezza nel verbo ``attendersi''. Da un lato, ci si
``attendeva'' (ossia si considerava probabile) che un Quarto cambiasse,
ma dall'altro il cambiamento da parte sua era anche ``atteso'' (ossia
preteso) dalla comunità e dai suoi pari.

Com'era cambiato Jason? Per prima cosa, si muoveva in modo diverso. Jase
aveva camuffato molto abilmente la sua SMA, ma nella camminata e nei
gesti c'era una nuova e percettibile disinvoltura. Era l'uomo di latta
dopo essere stato oliato. Di tanto in tanto era ancora di cattivo umore,
ma i suoi sbalzi erano meno violenti. Imprecava meno di frequente -
ovvero, era meno probabile che inciampasse in una di quelle voragini
emozionali in cui l'unico sostantivo utile è ``cazzo''. Scherzava più di
prima.

Tutte queste cose sembravano positive. E lo erano, ma erano anche
superficiali. Altri cambiamenti erano più preoccupanti. Si era ritirato
dalla gestione quotidiana della Perielio a un punto tale che il suo
staff lo informava una volta alla settimana e per il resto del tempo lo
ignorava. Aveva iniziato a leggere libri di astrofisica marziana a
partire dalle traduzioni originali, aggirando i protocolli di sicurezza
se non violandoli del tutto. L'unico evento che si era insinuato nella
sua ritrovata calma era la morte di Wun, e questo lo aveva lasciato
ferito e tormentato in un modo che ancora non capivo appieno.

«Ti rendi conto,» disse E.D., «che abbiamo appena visto la fine della
Perielio.»

E in un certo senso era così. A prescindere dall'interpretare qualunque
risposta noi ricevessimo dai replicatori, la Perielio in quanto agenzia
spaziale civile era finita. Il ridimensionamento era già iniziato
seriamente. Metà del personale di supporto era stato sospeso. I tecnici
stavano andando via più lentamente, attratti dalle università o dai
ricchi imprenditori.

«Allora così sia» ribatté Jason, ostentando o l'innato equilibrio di un
Quarto o un'ostilità a lungo soppressa per suo padre. «Abbiamo svolto il
lavoro che dovevamo fare.»

«Riesci a startene lì a dare questo verdetto? A me?»

«Credo sia vero.»

«Non importa che io abbia passato una vita a costruire ciò che tu hai
appena fatto a pezzi?»

«Importa?» Jason rifletté come se E.D. gli avesse realmente fatto una
domanda. «Fondamentalmente no, non credo.»

«Mio Dio, cosa ti è successo? Se commetti un errore di una tale
gravità\ldots»

«Non credo sia un errore.»

«\ldots{} dovresti assumertene la responsabilità.»

«Credo di averlo fatto.»

«Perché se fallisce, sarai l'unico a essere incolpato.»

«Lo capisco.»

«L'unico che faranno fuori.»

«Se sarà necessario.»

«Non posso proteggerti» disse E.D.

«Non hai mai potuto» rispose Jason.

\separatorefiori

Tornai alla Perielio con lui. In quei giorni Jason guidava un'auto
tedesca a celle a combustibile - un'auto di nicchia, dal momento che
quasi tutti avevamo ancora macchine a benzina progettate da persone che
non ritenevano ci fosse un futuro per cui valesse la pena preoccuparsi.
I pendolari ci superavano a tutta velocità nelle corsie di sorpasso,
affrettandosi verso casa prima che facesse buio.

Gli dissi che avevo intenzione di lasciare la Perielio e aprire un
ambulatorio per conto mio.

Jase rimase in silenzio per un po' guardando la strada, l'aria calda
faceva ribollire l'asfalto come se i confini del mondo si fossero
liquefatti per il calore. Poi disse: «Ma non devi, Tyler. La Perielio
dovrebbe tirare avanti ancora per qualche anno, e ho abbastanza potere
per continuare a tenerti a libro paga. Posso assumerti privatamente, se
è necessario.»

«Ma è questo il punto, Jase. Non \emph{è} necessario. Ero sempre un po'
sottoutilizzato alla Perielio.»

«Annoiato, vuoi dire?»

«Sarebbe bello sentirsi utile, per una volta.»

«Non ti senti utile? Se non fosse stato per te, sarei su una sedia a
rotelle.»

«Non è stato grazie a me. È stato grazie a Wun. Io ho solo spinto lo
stantuffo.»

«Niente affatto. Mi sei stato vicino durante il calvario. Lo apprezzo. E
poi\ldots{} ho bisogno di qualcuno con cui parlare, qualcuno che non
stia cercando di comprarmi o vendermi.»

«Quando è stata l'ultima volta che abbiamo avuto una vera
conversazione?»

«Solo perché ho superato una crisi medica non significa che non ce ne
sarà un'altra.»

«Sei un Quarto, Jase. Non avrai bisogno di un dottore per un'altra
cinquantina d'anni.»

«E le uniche persone che lo sanno siete tu e Carol. Ed è un altro motivo
per cui non voglio che te ne vada.» Esitò. «Perché non ti sottoponi al
trattamento anche tu? Ti concederesti altri cinquant'anni come minimo.»

Pensai di poterlo fare. Ma i cinquant'anni ci avrebbero portati in fondo
all'eliosfera del Sole in espansione. Sarebbe stato un gesto inutile.
«Preferirei essere utile adesso.»

«Sei assolutamente certo di volertene andare?»

E.D. avrebbe detto: ``Resta.''

E.D. avrebbe detto: ``È compito tuo prenderti cura di lui.''

E.D. avrebbe detto molte cose.

«Sì, assolutamente.»

Jason strinse il volante e fissò la strada come se vi avesse visto
qualcosa di infinitamente triste. «Bene» concluse. «Allora posso solo
augurarti buona fortuna.»

\separatorefiori

Il giorno in cui lasciai la Perielio, il personale di supporto mi
convocò in una delle sale del consiglio d'amministrazione ormai
utilizzate di rado per una festa d'addio, nella quale ricevetti il
genere di regali adatti per l'ennesimo allontanamento da una forza
lavoro in diminuzione: un cactus in miniatura in un vaso di terracotta,
una tazza da caffè con il mio nome sopra e una spilla da cravatta in
peltro a forma di caduceo.

Quella sera Jase si presentò alla mia porta con un regalo più
problematico.

Era una scatola di cartone legata con uno spago. Conteneva, quando
l'aprii, circa mezzo chilo di documenti stampati in maniera fitta e sei
dischi ottici non etichettati.

«Jase?»

«Informazioni mediche» disse. «Puoi considerarle come un libro di
testo.»

«Che tipo di informazioni mediche?»

Sorrise. «Dagli archivi.»

«Gli archivi marziani?»

Annuì.

«Ma sono informazioni segretate.»

«Sì, tecnicamente lo sono. Ma Lomax segreterebbe anche il numero di
telefono per il 911 se pensasse di poterla fare franca. Qui ci sono
informazioni che potrebbero mandare in rovina la Pfizer e la Eli Lilly.
Ma questa non è per me una preoccupazione legittima. E per te?»

«No, ma\ldots»

«Non penso nemmeno che Wun avrebbe voluto tenerle segrete. Così ho
distribuito, con discrezione, parti degli archivi qua e là, a persone di
cui mi fido. Non devi realmente farci qualcosa. Guardale o ignorale,
archiviale\ldots{} va bene comunque.»

«Fantastico. Grazie, Jase. Un regalo che potrebbe farmi arrestare.»

Il suo sorriso si allargò. «So che farai la cosa giusta.»

«Di qualunque cosa si tratti.»

«Lo scoprirai. Ho fiducia in te, Tyler. Fin dal momento della cura.»

«Cosa?»

«Mi sembra di vedere le cose più chiaramente» disse.

Non diede spiegazioni e misi la scatola nella valigia come una specie di
souvenir. Fui tentato di scriverci sopra la parola {RICORDI}.

\separatorefiori

La tecnologia dei replicatori era lenta perfino a confronto con la
terraformazione di un pianeta morto. Passarono due anni prima che i
carichi utili che avevamo disseminato tra i planetesimi ai confini del
sistema solare rilevassero qualcosa di simile a un responso.

I replicatori là fuori erano occupati, sebbene a malapena sfiorati dalla
gravità del Sole, a fare quello per cui erano stati progettati:
riprodursi poco a poco e secolo dopo secolo, seguendo le istruzioni
scritte nel loro equivalente superconduttivo del DNA. Con il tempo e
un'adeguata riserva di ghiaccio ed elementi con tracce di carbonio, alla
fine avrebbero chiamato casa. Ma i primi, pochi satelliti rilevatori
posizionati oltre la membrana dello Spin ricaddero sulla Terra senza
registrare alcun segnale.

Durante quei due anni riuscii a trovare un socio (Herbert Hakkim, un
cordiale dottore originario del Bengali che aveva concluso il suo
internato l'anno in cui Wun visitò il Grand Canyon), e insieme
subentrammo a un medico generico che stava andando in pensione in un
ambulatorio di San Diego. Hakkim era aperto e cordiale con i pazienti ma
non aveva una vera vita sociale e sembrava preferire così: ci trovavamo
di rado al di fuori dell'orario di lavoro, e credo che la domanda più
intima che mi abbia mai posto fosse perché avevo due telefoni cellulari.
(Uno per motivi di abitudine; l'altro perché il numero assegnatogli era
l'ultimo che avevo dato a Diane. Non che avesse mai squillato. Né io
avevo tentato di ricontattarla. Ma se avessi lasciato scadere il numero,
lei non avrebbe avuto modo di mettersi in contatto con me, e questo mi
sembrava ancora\ldots{} be', sbagliato.)

Mi piaceva il mio lavoro, e in generale mi piacevano i miei pazienti. Ho
visto più ferite d'arma da fuoco di quanto un tempo mi sarei aspettato,
ma quelli erano gli anni duri dello Spin; le curve di tendenza nazionali
degli omicidi e dei suicidi avevano cominciato ad andare in direzione
verticale. Gli anni in cui sembrava che chiunque al di sotto dei
trent'anni indossasse una qualche uniforme: le forze armate, la Guardia
Nazionale, il Dipartimento per la Sicurezza Nazionale, la sicurezza
privata; perfino gli scout domestici e le guide coniugali per far fronte
al calo della natalità. Gli anni in cui Hollywood iniziò a sfornare film
ultraviolenti o ultrareligiosi nei quali, comunque, lo Spin non era mai
menzionato in maniera esplicita; lo Spin, come il sesso e le parole che
lo descrivevano, era stato bandito dalla ``dialettica
dell'intrattenimento'' da parte del Comitato culturale di Lomax e dalla
Commissione federale per le comunicazioni.

Quelli erano anche gli anni in cui l'amministrazione mise in atto un
mucchio di nuove leggi mirate all'epurazione degli archivi marziani. Gli
archivi di Wun, secondo il Presidente e i suoi alleati del Congresso,
contenevano conoscenze intrinsecamente pericolose che dovevano essere
revisionate e messe al sicuro. Renderli accessibili al pubblico sarebbe
stato come ``pubblicare in Internet i piani di un ordigno nucleare
portatile''. Perfino il materiale antropologico fu vagliato
attentamente: nella versione pubblicata un Quarto era definito come ``un
anziano rispettabile''. Nessuna menzione della longevità ottenuta
mediante l'intervento medico.

Ma chi aveva bisogno o voleva la longevità? La fine del mondo era ogni
giorno più vicina.

Gli sfarfallii ne erano la prova, se c'era qualcuno che aveva bisogno di
prove.

\separatorefiori

Quando cominciarono gli sfarfallii, i primi risultati positivi del
progetto dei replicatori erano giunti da sei mesi.

Appresi la maggior parte delle notizie sui replicatori da Jase un paio
di giorni prima che queste facessero irruzione nei media. Di per sé, non
erano niente di spettacolare. Un satellite di rilevamento della NASA /
Perielio aveva registrato un debole segnale proveniente da un corpo
conosciuto nella Nube di Oort ben oltre l'orbita di Plutone - un bip
intermittente non codificato che era il suono di una colonia di
replicatori prossima al completamento. (Si potrebbe dire, prossima alla
maturazione.)

Sarebbe potuta sembrare una cosa di poco conto se non se ne fosse
considerato il significato.

Le cellule dormienti di una biologia artificiale del tutto nuova si
erano imbattute in un grosso pezzo di ghiaccio polveroso nello spazio
più profondo. Quelle cellule avevano, in seguito, dato inizio a una
forma atrocemente lenta di metabolismo, nel quale assorbivano lo scarso
calore del Sole distante, lo usavano per separare alcune molecole vicine
di acqua e carbonio, e si duplicavano con le materie prime che ne
derivavano.

Nel corso di molti anni la stessa colonia crebbe fin quasi a raggiungere
le dimensioni di un cuscinetto a sfera. Un astronauta che avesse fatto
quel lunghissimo viaggio e saputo con esattezza dove guardare, l'avrebbe
vista come una fossetta scura sulla regolite rocciosa e ghiacciata del
planetesimo ospite. Ma la colonia a livello infinitesimale era più
efficiente del suo antenato unicellulare. Cominciò a crescere più
velocemente e generare maggior calore. Il differenziale di temperatura
tra colonia e ambiente circostante era solo di una frazione di grado
Kelvin (tranne quando brevi esplosioni riproduttive pompavano energia
latente nell'ambiente locale), ma era costante.

Trascorsero molti altri millenni (o mesi terrestri.) Le subroutine nel
sostrato genetico dei replicatori, attivate dai gradienti termici
locali, alterarono la proliferazione della colonia. Le cellule
cominciarono a differenziarsi. Come un embrione umano, la colonia non si
limitò a produrre più cellule ma cellule specializzate, l'equivalente di
cuore e polmoni, braccia e gambe. Alcuni filamenti si spinsero fin
dentro il materiale poco compatto del planetesimo, scavandolo per
estrarre molecole carboniose.

Infine, esplosioni di vapore microscopiche ma attentamente calcolate
iniziarono a rallentare la rotazione dell'oggetto ospite (in modo
paziente, nel corso dei secoli), finché la superficie della colonia si
trovò rivolta perennemente verso il Sole. A quel punto la
differenziazione cominciò sul serio. La colonia estruse legami di
carbonio-carbonio e carbonio-silicio; sviluppò filamenti monomolecolari
per congiungere quei legami, evolvendosi nella scala della complessità;
da quei legami generò punti fotosensibili - occhi - e la capacità di
generare e riprodurre microimpulsi di suono a radiofrequenza.

Trascorsi altri secoli, la colonia elaborò e affinò quelle capacità
finché si annunciò con un semplice segnale acustico intermittente,
l'equivalente del verso di un passero appena nato. Che era ciò che i
nostri satelliti avevano rilevato.

Per qualche giorno i mezzi d'informazione pubblicarono la storia (con
materiale d'archivio di Wun Ngo Wen, del suo funerale e del lancio) e
poi se ne dimenticarono. Dopotutto, quella era solo la prima fase di ciò
che i replicatori dovevano fare secondo la loro progettazione.

Niente di più. Poco interessante. A meno di non pensarci per più di
trenta secondi.

Si trattava letteralmente di tecnologia con una vita propria. Un genio
fuori dalla lampada per sempre.

\separatorefiori

Gli scintillii si manifestarono qualche mese più tardi.

Gli scintillii furono il primo segnale di cambiamento o interferenza
nella membrana dello Spin\ldots{} il primo, diciamo, a meno di tener
conto dell'evento che seguì l'attacco missilistico cinese sugli
artefatti polari, nei primissimi anni dello Spin. Entrambi gli
avvenimenti furono visibili da qualunque punto del pianeta. Ma a parte
quella somiglianza fondamentale, non avevano altri elementi in comune.

Dopo l'attacco missilistico, la membrana dello Spin sembrò funzionare a
singhiozzo e riprendersi, generando immagini stroboscopiche del cielo in
evoluzione, molteplici lune e stelle che roteavano.

Gli scintillii erano diversi.

Li osservai dalla terrazza del mio appartamento in periferia. Una calda
sera di settembre. Alcuni vicini si trovavano già all'aperto quando gli
scintillii ebbero inizio. Poi uscimmo tutti. Ci appollaiammo sui
parapetti come storni, schiamazzando.

Il cielo s'illuminò.

Non per via delle stelle ma per via dei fili infinitesimamente sottili
di fiamme dorate, che crepitavano come folgori senza calore da un
orizzonte all'altro. I fili si spostavano e mutavano in modo
imprevedibile; alcuni apparivano o scomparivano del tutto; di tanto in
tanto altri si manifestavano in una vampata di fuoco. Era affascinante
ma allo stesso tempo terrificante.

L'avvenimento fu globale, non locale. Sulla parte illuminata del
pianeta, il fenomeno fu visibile solo lievemente, attenuato dalla luce
solare o coperto dalle nuvole; nel Nord e nel Sud America e in Europa
occidentale lo spettacolo del cielo scuro causò sporadiche ondate di
panico. Dopotutto, avevamo atteso la fine del mondo da molti più anni di
quanti ne volessimo ricordare. Sembrava, quantomeno, il preludio alla
realtà.

Quella notte si verificarono centinaia di suicidi o tentativi di
suicidio, e in più, decine di omicidi o eutanasie nella città dove
vivevo. In tutto il mondo, le cifre furono immensamente più alte. A
quanto pare, ci furono moltissime persone come Molly Seagram che
scelsero di evitarsi la tanto preannunciata ebollizione degli oceani con
qualche pasticca letale di questo o di quello. Farmaci mortali tenuti da
parte fino a quel momento per sé e per familiari e amici. Molti di loro
optarono per la via di fuga definitiva non appena il cielo s'illuminò.
Una scelta prematura, per come andò a finire.

Lo spettacolo durò otto ore. Al mattino mi trovavo all'ospedale locale a
dare una mano nel reparto di pronto soccorso. Per mezzogiorno avevo
visto sette diversi casi di avvelenamento da monossido di carbonio,
gente che si era intenzionalmente rinchiusa nel proprio garage con
l'auto in folle. Quasi tutti erano morti ben prima che ne dichiarassi il
decesso e i sopravvissuti stavano poco meglio di loro. Persone che
altrimenti avrebbero potuto star bene, persone che avrei potuto
incontrare dal droghiere, avrebbero trascorso il resto della loro vita
attaccate a dei respiratori, con irreparabili danni al cervello, vittime
di una strategia di fuga maldestra. Per niente piacevole. Ma le ferite
d'arma da fuoco alla testa erano peggio. Non riuscivo a occuparmi di
loro senza pensare a Wun Ngo Wen steso su quell'autostrada in Florida,
col sangue che gocciolava da ciò che restava del suo cranio.

Otto ore. Dopodiché il cielo tornò a rischiararsi e il Sole ne uscì
splendente come la battuta finale di una brutta barzelletta.

Accadde nuovamente un anno e mezzo più tardi.

\separatorefiori

«Hai l'aria di un uomo che ha perso la fede» mi disse una volta Hakkim.

«O che non ne ha mai avuta una» risposi.

«Non intendo fede in Dio. Su quel fronte sembri essere sinceramente
innocente. Fede in qualcos'altro. Non so cosa.»

Il che sembrò criptico. Ma lo compresi un po' meglio quando parlai con
Jason la volta seguente.

Ero a casa quando mi chiamò. (Al mio solito numero di cellulare, non sul
telefono orfano che mi portavo dietro come un amuleto sfortunato.)
Risposi: «Pronto.» E lui: «Devi assolutamente vedere questa cosa alla
televisione.»

«Vedere cosa?»

«Sintonizzati su un notiziario. Sei solo?»

La risposta era sì. Per mia scelta. Nessuna Molly Seagram a complicarmi
la fine dei tempi. Il telecomando della TV era sul tavolino da caffè
dove lo avevo lasciato. Dove lo lasciavo sempre.

Il canale delle notizie mostrava un grafico molto colorato accompagnato
da una blaterante voce fuoricampo. Spensi l'audio. «Cosa sto guardando,
Jase?»

«Una conferenza stampa del JPL. La serie di dati recuperata dall'ultimo
satellite ricevitore orbitale.»

In altre parole, i dati dei replicatori. «E quindi?»

«Siamo in affari» disse, e potevo praticamente sentirne il sorriso.

Il satellite aveva rilevato molteplici sorgenti radio che trasmettevano
dall'esterno del sistema solare. Il che significava che più di una
colonia di replicatori era giunta alla maturazione. ``E i dati erano
complessi,'' spiegò Jason, ``non semplici.'' Mentre le colonie di
replicatori invecchiavano, il loro ritmo di crescita rallentava ma le
loro funzioni diventavano più affinate e mirate. Non si stavano più
semplicemente rivolgendo verso il Sole in cerca di energia gratuita.
Stavano analizzando la luce delle stelle, calcolando le orbite dei
pianeti su reti neurali fatte di silicio e fibre di carbonio, mettendole
a confronto con i modelli impressi nel loro codice genetico. Almeno una
decina di colonie completamente adulte avevano rimandato esattamente i
dati che erano state progettate per raccogliere, quattro flussi di dati
binari che dichiaravano:

1. Quello era il sistema planetario di una stella con massa solare 1.0;

2. Il sistema possedeva otto grandi corpi celesti (Plutone era sceso al
di sotto del limite di massa rilevabile);

3. Due di quei pianeti erano invisibili, circondati dalle membrane dello
Spin; e

4. Le colonie di replicatori trasmittenti erano passate alla modalità
riproduttiva, liberandosi di cellule-seme non specifiche e lanciandole
con esplosioni di vapore cometario verso le stelle vicine.

Lo stesso messaggio, disse Jason, era stato trasmesso alle colonie
locali e meno mature, che avrebbero reagito escludendo le funzioni
superflue e rivolgendo la loro energia verso un comportamento puramente
riproduttivo.

In altre parole, avevamo infettato con successo il sistema solare
esterno con i sistemi semibiologici di Wun.

Che adesso erano in fase di sporulazione.

«Questo non ci dirà niente dello Spin» obiettai.

«Certo che no. Non ancora. Ma questo piccolo rigagnolo di informazioni
diventerà presto un torrente. Con il tempo saremo in grado di elaborare
una mappa dello Spin di tutte le stelle vicine\ldots{} e alla fine,
forse, dell'intera galassia. Da essa dovremmo poter dedurre da dove
vengono gli Ipotetici, dove hanno posizionato gli altri Spin e cosa
accade alla fine ai mondi avvolti dallo Spin quando le loro stelle si
espandono e si estinguono.»

«Comunque questo non sistemerà niente, vero?»

Sospirò come se lo avessi deluso per aver fatto una domanda stupida.
«Probabilmente no. Ma non è meglio sapere piuttosto che fare congetture?
Potremmo scoprire di essere condannati, o potremmo scoprire di avere più
tempo di quanto pensavamo. Ricorda, Tyler, che stiamo lavorando anche su
altri fronti. Stiamo scavando a fondo nella fisica teorica degli archivi
di Wun. Se si modella la membrana dello Spin come un tunnel
spazio-temporale, il quale avvolge un oggetto che accelera quasi alla
velocità della luce\ldots»

«Ma noi non stiamo accelerando. Non stiamo andando da nessuna parte.» Se
non a capofitto verso il futuro.

«No, ma se fai il calcolo, vedrai che i risultati combaciano con il
nostro studio dello Spin. Il che potrebbe darci un'indicazione in merito
a \emph{quali forze} stanno manipolando gli Ipotetici.»

«A quale scopo, però, Jase?»

«Troppo presto per dirlo. Ma non credo nell'inutilità della conoscenza.»

«Anche se stiamo morendo?»

«Tutti muoiono.»

«Voglio dire, come specie.»

«Questo è da vedere. Qualunque cosa sia lo Spin, dev'essere più di una
sorta di eutanasia globale elaborata. Gli Ipotetici devono avere
sicuramente uno scopo.»

Forse era così. Ma mi resi conto che quella era la fede che mi aveva
abbandonato. La fede nella Grande Salvezza.

La Grande Salvezza, un marchio di prodotti per tutti i gusti. All'ultimo
minuto avremmo potuto ideare una soluzione tecnologica e salvarci.
Oppure: gli Ipotetici erano creature benevole che volevano trasformare
il pianeta in un regno pacifico. Oppure: Dio avrebbe potuto salvarci
tutti, o almeno i veri credenti tra di noi. Oppure. Oppure. Oppure.

La Grande Salvezza. Una bugia stucchevole. Una scialuppa di salvataggio
di carta, anche se ci stavamo uccidendo per cercare di aggrapparvici.
Non era lo Spin ad aver mutilato la mia generazione. Era il miraggio e
la ricompensa della Grande Salvezza.

\separatorefiori

Gli scintillii ricomparvero l'inverno seguente, durarono quarantaquattro
ore, poi scomparvero di nuovo. Molti di noi iniziarono a considerarli
una specie di fenomeno atmosferico celeste, imprevedibile ma
generalmente innocuo.

I pessimisti fecero notare che gli intervalli tra gli episodi
diventavano sempre più brevi, e la loro durata sempre più lunga.

Ad aprile si verificarono degli scintillii che durarono tre giorni e che
interferirono con la trasmissione dei segnali dell'aerostato. Questi
provocarono un'altra ondata (più piccola) di suicidi e tentati suicidi -
persone in preda al panico più per il malfunzionamento dei loro telefoni
e televisori che per ciò che vedevano in cielo.

Avevo smesso di prestare attenzione alle notizie, ma alcuni eventi erano
impossibili da ignorare: le sconfitte militari in Nordafrica e in Europa
dell'Est, il colpo di Stato da parte di una setta in Zimbabwe, i suicidi
di massa in Corea. Quell'anno i rappresentanti dell'Islam apocalittico
ottennero grandi risultati nelle elezioni in Algeria e in Egitto. Un
culto filippino che venerava il ricordo di Wun Ngo Wen - da loro eletto
a santo pastorale, un Gandhi agricolo - era riuscito a organizzare uno
sciopero generale a Manila.

E io ricevetti qualche chiamata in più da parte di Jason. Mi inviò per
posta un telefono con una sorta di sistema di cifratura integrato, il
quale, affermò, ci avrebbe fornito ``un'ottima difesa contro i
cacciatori di parole chiave'', qualunque cosa volesse dire.

«Sembra un po' paranoico» riflettei.

«Paranoico in modo utile, penso.»

Forse, se avessimo voluto discutere di questioni riguardanti la
sicurezza nazionale. Tuttavia, non lo facemmo, almeno non all'inizio.
Invece Jason mi chiese del mio lavoro, della mia vita e della musica che
ascoltavo. Capii che stava cercando di dare vita a quel genere di
conversazione che avremmo potuto avere venti o trent'anni prima\ldots{}
prima della Perielio, se non prima dello Spin. Mi raccontò che era stato
a trovare sua madre. Carol continuava a scandire i suoi giorni con
orologio e bottiglia. Niente era cambiato. Su quel punto, Carol era
stata irremovibile. Il personale della casa teneva tutto pulito, tutto
al suo posto. La Grande Casa era come una macchina del tempo, disse
Jason, come se fosse stata sigillata ermeticamente la prima sera dello
Spin. Era un po' inquietante.

Domandai se Diane avesse più chiamato.

«Diane ha smesso di rispondere a Carol da prima che Wun venisse ucciso.
No, non una parola da parte sua.»

Allora gli chiesi del progetto dei replicatori. Negli ultimi tempi non
era apparso niente sui giornali.

«Non disturbarti a cercare. La JPL sta insabbiando i risultati.»

Sentii la tristezza nella sua voce. «È così grave?»

«Non sono notizie del tutto brutte. Almeno, non fino a ora. I
replicatori hanno fatto tutto ciò che Wun sperava facessero. Cose
fantastiche, Tyler. Intendo, assolutamente fantastiche. Spero di poterti
mostrare le mappe che abbiamo realizzato. Grandi mappe sotto forma di
software esplorabili. Quasi duecentomila stelle, in un alone di
centinaia di anni luce di diametro. Adesso sappiamo sull'evoluzione
delle stelle e dei pianeti più di quanto un astronomo della generazione
di E.D. avrebbe potuto immaginare.»

«Ma niente riguardo allo Spin?»

«Non ho detto questo.»

«Allora cosa avete appreso?»

«Innanzitutto, che non siamo soli. In quel volume di spazio abbiamo
trovato tre pianeti invisibili approssimativamente delle dimensioni
della Terra, in orbite che secondo i parametri terrestri sono abitabili,
o lo sono state in passato. Il più vicino orbita attorno alla stella 47
Ursae Majoris. Il più lontano\ldots»

«Non ho bisogno dei dettagli.»

«Se prendiamo in considerazione l'età delle stelle interessate e
formuliamo delle ipotesi plausibili, sembrerebbe che gli Ipotetici
provengano da qualche parte in direzione del centro della galassia. Ci
sono anche altri indizi. I replicatori hanno trovato una coppia di
stelle nane bianche - stelle estinte, fondamentalmente, ma che avevano
un aspetto simile a quello del Sole qualche miliardo di anni fa - con
pianeti rocciosi in orbita che non devono essere sopravvissuti
all'espansione solare.»

«Superstiti dello Spin?»

«Forse.»

«Sono pianeti \emph{vivi}, Jase?»

«Non abbiamo modo di saperlo. Ma non hanno membrane dello Spin che li
proteggano, e il loro attuale habitat stellare è del tutto ostile per i
nostri standard.»

«Che sarebbe a dire?»

«Non lo so. Nessuno lo sa. Pensavamo che saremmo stati in grado di fare
dei raffronti più significativi non appena la rete di replicatori si
fosse estesa. Ciò che abbiamo creato con i replicatori è veramente una
rete neurale su scala vasta a livelli inimmaginabili. I replicatori
comunicano tra di loro nello stesso modo in cui comunicano i neuroni, la
differenza è che loro lo fanno attraverso i secoli e gli anni luce.
Quello che fanno è incredibilmente affascinante, davvero. Una rete più
grande di qualsiasi cosa il genere umano abbia mai costruito. Che
raccoglie informazioni, le seleziona, le immagazzina, le ritrasmette a
noi\ldots»

«Allora cos'è andato storto?»

Sembrò come se soffrisse nel dirlo. «Forse l'invecchiamento. Tutto
invecchia, anche i codici genetici altamente protetti. Forse si stanno
evolvendo oltre le nostre istruzioni. Oppure\ldots»

«Sì, ma cos'è \emph{successo}, Jase?»

«I dati stanno diminuendo. Stiamo ricevendo informazioni frammentarie e
contraddittorie dai replicatori che si trovano più distanti dalla Terra.
Questo potrebbe voler dire molte cose. Se stanno morendo, ciò potrebbe
riflettere il delinearsi di alcuni errori nel codice di progettazione.
Ma anche alcuni punti di ricezione consolidati da tempo stanno iniziando
a disattivarsi.»

«Qualcosa li sta prendendo di mira?»

«È un'ipotesi troppo affrettata. Ecco un'altra idea. Quando abbiamo
lanciato quelle cose verso la Nube di Oort, abbiamo dato vita a
un'ecologia interstellare semplice: ghiaccio, polvere e vita
artificiale. Ma se non fossimo stati i primi? Se l'ecologia
interstellare non fosse semplice?»

«Vuoi dire che potrebbero esserci altri tipi di replicatori là fuori?»

«Può darsi. Se così fosse, potrebbero essere in competizione per
ottenere risorse. Potrebbero addirittura servirsi gli uni degli altri
per ottenere risorse. Pensavamo di aver inviato i nostri dispositivi in
un vuoto sterile. Ma potrebbero esserci specie rivali, potrebbero
esserci perfino specie predatrici.»

«Jason\ldots{} pensi che qualcosa li stia \emph{divorando}!»

«È possibile» rispose.

\separatorefiori

Gli scintillii si manifestarono di nuovo nel mese di giugno e durarono
quasi quarantotto ore prima di dissolversi.

Ad agosto, cinquantasei ore di scintillii e in aggiunta problemi
intermittenti con le telecomunicazioni.

Quando iniziarono nuovamente verso la fine di settembre, nessuno ne
rimase sorpreso. Trascorsi quasi tutta la prima sera con le persiane
chiuse, ignorando il cielo e guardando un film che avevo scaricato una
settimana prima. Un vecchio film, antecedente allo Spin. Lo guardai non
per la trama ma per i volti, i volti delle persone com'erano un tempo,
persone che non avevano trascorso le loro vite temendo per il futuro.
Persone che, di tanto in tanto, parlavano della Luna e delle stelle
senza ironia o nostalgia.

A quel punto squillò un cellulare.

Non il mio telefono personale, e nemmeno quello criptato che Jase mi
aveva spedito. Riconobbi immediatamente il suono a tre toni sebbene non
lo sentissi da anni. Era percepibile ma flebile\ldots{} flebile perché
lo avevo lasciato nella tasca di una giacca appesa nell'armadio del
corridoio.

Squillò altre due volte prima che lo tirassi fuori in modo maldestro e
dicessi: «Pronto?»

Mi aspettavo che qualcuno avesse sbagliato numero. Desideravo sentire la
voce di Diane. Lo desideravo e lo temevo.

Ma all'altro capo c'era una voce da uomo. Simon. Lo riconobbi in
ritardo.

Chiese: «Tyler? Tyler Dupree? Sei tu?»

Avevo ricevuto tante chiamate d'emergenza da accorgermi dell'ansia nella
sua voce. Risposi: «Sono io, Simon. Cosa c'è che non va?»

«Non dovrei parlare con te. Ma non so chi altro chiamare. Non conosco
alcun medico del posto. E lei sta molto male. Sta davvero molto male,
Tyler. Non sembra che migliori. Credo che abbia bisogno di\ldots»

E a quel punto gli scintillii interruppero la comunicazione e non si udì
altro che un crepitio.